\section{DISCUSSION AND 6G RESEARCH ROADMAP}
\label{sec:discussion_roadmap}

\subsection{A Cross-Modality Field Rather Than a Single Design Family}

The 206-study corpus shows that O-ISAC has developed as a family of related
optical and photonic paradigms rather than as one homogeneous system class.
Photonic-terahertz and fiber platforms together contributed 125 studies, but
VLC/LiFi and FSO added 69 further studies with markedly different propagation,
hardware, illumination, mobility, and sensing constraints. Hybrid and other
optical platforms formed a smaller frontier. A useful O-ISAC taxonomy must
therefore preserve the physical medium and sensing construct instead of
projecting all studies onto an RF-inspired joint-radar-and-communication
template.

At the same time, the integration results reveal a common design grammar.
Shared resource allocation, shared hardware, and shared waveforms each
appeared in more than half of the corpus, while shared links/channels and joint
optimization were also prominent. This commonality provides a credible basis
for a unified survey, provided that ``integration'' is reported at the level
actually demonstrated. Reuse of a carrier, waveform, front end, link, sensing
observable, or application context are different architectural commitments
and should not be collapsed into a binary ISAC label.

\subsection{Metric Abundance Does Not Equal Comparability}

The evidence model contains 4,779 primary metric claims, spanning sensing,
communication, joint, implementation, system, energy, security,
machine-learning, and other domains. This richness is scientifically useful,
but it also explains why a universal performance ranking would be misleading.
Metric names alone rarely define comparable evidence. A numerical value is
interpretable only together with its metric role, measurement plane, operating
condition, scenario, and validation context.

The appraisal makes this distinction visible. Reporting completeness was
strong in 196 studies and metric clarity was strong in 168, yet only 4 studies
received the strongest comparison-admissibility score and none reached the
strongest benchmark-readiness score. Many reports are therefore sufficiently
clear to understand on their own terms but insufficiently standardized for
direct cross-study ranking. Future O-ISAC reporting should attach a compact
metric contract to each headline result: definition, numerator and
denominator, physical or electrical plane, reference bandwidth or time,
scenario, channel/target model, hardware state, and uncertainty or repeat
count.

\subsection{Trade-offs Are Multi-Resource and Task-Dependent}

The 404 trade-off records do not support reducing O-ISAC design to a single
rate--resolution curve. Bandwidth, spectrum, or resource allocation was the
largest trade-off family, followed closely by power, energy, or dynamic range.
Communication reliability versus sensing quality, rate--resolution,
rate--accuracy, and rate--localization were also substantial. Range, coverage,
hardware complexity, and security formed additional axes.

This structure suggests that practical O-ISAC optimization should expose at
least four coupled budgets: spectral/time resources, optical/electrical power
and dynamic range, hardware/DSP complexity, and task-specific performance.
Pareto comparisons are meaningful only when these budgets and the validation
plane are aligned. A curve obtained from a simulation under one channel model
should not be presented as directly dominating a laboratory or field result
that uses a different sensing observable or performance definition.

\subsection{Experimental Momentum, Reproducibility Deficit}

Three-quarters of the corpus reached at least a laboratory experiment,
proof-of-concept, controlled prototype, or field tier, demonstrating that
O-ISAC is not merely a theoretical research program. The maturity profile is
nevertheless uneven: only 12 studies reached field or deployment-oriented
validation. More importantly, experimental maturity was not accompanied by
open research artifacts. Only 13 studies provided open data, while code or
model artifacts were unavailable or not reported in 197.

This gap is a central barrier to cumulative progress. A convincing prototype
can establish feasibility, but without reusable waveforms, channel/target
traces, calibration procedures, parameter files, code, or baseline outputs, it
cannot easily become a community benchmark. The near-term priority is not one
universal dataset for all optical modalities. It is a family of modality-aware
benchmark packages with a common metadata envelope and explicit translation
rules between measurement planes.

\subsection{Enabling Technologies and Application Pull}

The multi-label technology map shows several overlapping streams. These
include photonic-terahertz and microwave-photonic generation; FMCW and other
chirped processing; coherent optics; multicarrier modulation; distributed
fiber sensing; integrated photonics; beam control; and learning-based
inference. No single enabler defines O-ISAC. Instead, the field advances by
recombining optical generation, propagation, detection, waveform, and
inference technologies around a shared sensing--communication objective.

The application map is similarly broad. 6G/access-network contexts were most
frequent, followed by environmental monitoring, vehicular systems, smart
infrastructure, industrial settings, optical access networks, security, and
indoor positioning. The 6G-relevance crosswalk classified 138 studies as
directly relevant and 64 as inferentially relevant. The evidence therefore
supports O-ISAC as a substantive component of the 6G research landscape, but
the label should remain tied to demonstrated capabilities rather than used as
a substitute for system-level requirements or deployment evidence.

\subsection{Research Roadmap}

The synthesis supports six priorities.

\begin{enumerate}
\item \emph{Adopt measurement-plane contracts.} Headline sensing and
communication values should identify their physical/electrical reference
plane, role, bandwidth/time basis, and uncertainty.
\item \emph{Publish modality-aware benchmark bundles.} Data, code or models,
waveforms, calibration records, parameter files, and baseline outputs should
be released together under a common metadata structure.
\item \emph{Move from feasibility to field evidence.} Laboratory prototypes
should be tested under mobility, environmental variation, component drift,
interference, and realistic network load.
\item \emph{Report multi-resource Pareto fronts.} Evaluations should expose
spectral, power, complexity, and task-performance budgets instead of selecting
one favorable trade-off axis.
\item \emph{Separate shared-component claims from joint optimization.}
Architecture reporting should identify what is shared, where the functions
separate, and whether sensing information closes a communication-control loop
or vice versa.
\item \emph{Design cross-modality comparisons around equivalent tasks.}
Benchmarking should compare matched operational objectives and constraints,
not nominally similar metrics that represent different sensing constructs.
\end{enumerate}

\subsection{Limitations of the Evidence and Review Process}

The review is limited by heterogeneity in terminology, measurement planes,
scenarios, and validation practices. The absence of common effect measures
precluded conventional meta-analysis, and the descriptive distributions
should not be interpreted as pooled causal effects. Fifty-eight unique reports
were not retrieved. Contextual records were kept outside the primary technical
corpus, while 31 context-only claims and 72 claims with unresolved source
conflicts were excluded from primary numerical synthesis. Claim-level
quarantine preserves uncertainty but can reduce the evidence available for a
specific comparison without excluding the entire study.

The workflow was investigator-supervised, AI-assisted, user-delegated, and
claim-governed. Independent duplicate human screening and independent
full-corpus human PDF verification are not documented; inter-rater reliability
is therefore not reported. Deterministic QA, source locators, decision logs,
hashes, and explicit conflict registers improve auditability but do not replace
reviewer independence. These limitations should be considered when
interpreting fine-grained prevalence estimates or claims that depend on a
small evidence body.

\section{CONCLUSION}
\label{sec:conclusion_206}

This review consolidates 227 eligible reports into 206 unique O-ISAC studies
and organizes 8,203 primary claims through a cross-modality, metric-governed
synthesis. The field has established a broad technical base spanning
photonic-terahertz, fiber, VLC/LiFi, FSO, and hybrid optical platforms, with
resource, hardware, waveform, and link reuse providing a shared integration
grammar. Its main limitation is no longer a shortage of reported performance
values; it is the lack of portable measurement contracts, common baselines,
open artifacts, and field evidence required to compare those values
responsibly. Progress toward 6G-ready O-ISAC will therefore depend on coupling
continued photonic and signal-processing innovation with reproducible,
task-equivalent, and modality-aware evaluation. The taxonomy, claim-governance
rules, technical appraisal, and S1--S7 evidence maps developed here provide an
auditable foundation for that transition.
